% ====================================================================
% §2 통계역학 기초 (Part 0) — 파트 a: 앙상블·단일자리·격자기체 (2.1–2.3)
% ====================================================================
\section{통계역학 기초 --- 분배함수에서 Nernst 식까지 (Part 0)}\label{sec:sm-found}

본론($\mathrm N1$--$\mathrm N9$)의 모든 평형식은 통계역학의 한 사슬에서 나온다. 이 절은 그 사슬을 처음부터 놓으며, 통계역학
사전지식 없이도 본론의 출발점들이 이 절 안에서 닫히도록 각 단계를 (a) 출발 $\to$ (b) 연산 $\to$ (c) 중간식
$\to$ (d) 박스로 전개한다. 유도의 골격은 고전 통계역학(분배함수$\cdot$화학퍼텐셜)으로 닫히고, 양자 통계가
필요한 항(점유 자리의 진동 자유도, Chapter 2 의 포논$\cdot$전자 분포)은 해당 자리에서 명시적으로 도입한다.
사슬은
\[
\boxed{\;\begin{aligned}
&\text{미시상태·앙상블}\ \to\ Z,\Xi\ \to\ \theta=\frac{1}{1+e^{\beta(\tilde\varepsilon-\mu)}}\ \to\ \text{lattice gas}\\
&\to\ \mu(\theta;\Omega)\ \to\ \xi_\eq\ \text{logistic}\ \to\ G,\ \text{Nernst}
\end{aligned}\;}
\]
이고, 도착점들이 본론 결과 사슬 --- \S\ref{sec:center} 의 $U_j(T)$, \S\ref{sec:hys} 의 히스테리시스,
\S\ref{sec:width} 의 $\xi_\eq$ --- 의 입력이다.

\subsection{앙상블과 Boltzmann·Gibbs 인자 --- 분배함수와 화학퍼텐셜 $\mu$}\label{sec:sm-ensemble}
\textbf{(a) 출발 --- 미시상태·앙상블과 두 공준.} \emph{미시상태}(microstate)는 계의 모든 자유도를 지정한
배치 하나이고, \emph{거시상태}(macrostate)는 에너지·부피·입자수 $(E,V,N)$ 같은 소수의 거시변수만 지정한
미시상태의 집합이다. 같은 거시 조건을 공유하는 미시상태들의 모음이 \emph{앙상블}(ensemble)이다. 통계역학은
두 공준으로 시작한다 --- 고립계 평형에서 접근 가능한 $W$ 개 미시상태는 등확률이고(공준 1), 엔트로피는 그
경우의 수의 로그다(공준 2, Boltzmann):
\begin{equation}
S\;=\;k_B\ln W.
\label{eq:sm-S}
\end{equation}
열역학과의 접점은 근본 미분식 $\dd E=T\dd S-P\dd V+\mu\,\dd N$ 을 $S$ 에 대해 푼 형태다:
\begin{equation}
\dd S=\frac{1}{T}\,\dd E+\frac{P}{T}\,\dd V-\frac{\mu}{T}\,\dd N
\quad\Longrightarrow\quad
\frac{\partial S}{\partial E}\Big|_{V,N}=\frac1T,\qquad
\frac{\partial S}{\partial N}\Big|_{E,V}=-\frac{\mu}{T}.
\label{eq:sm-fund}
\end{equation}
여기서 $\mu$ 의 물리적 의미가 이미 보인다 --- 입자 1개를 받아들일 때 엔트로피가 $\mu/T$ 만큼 \emph{감소}하는
값비쌈의 척도다. 두 계 1,2 가 입자를 교환하면 총 엔트로피 변화가
$\dd S_\mathrm{tot}=\big(\tfrac{\mu_2}{T}-\tfrac{\mu_1}{T}\big)\dd N_1\ge0$ 이므로 입자는 $\mu$ 가 높은 쪽에서
낮은 쪽으로 흐르고, 정지 조건은 $\mu_1=\mu_2$ --- 화학퍼텐셜의 일치가 곧 입자 교환 평형이다.

\textbf{(b) 연산 --- 저장조와 접촉한 작은 계.} 관심 계가 거대한 저장조(reservoir)와 에너지·입자를 교환한다.
계가 미시상태 $i$(에너지 $E_i$, 입자수 $N_i$)에 있을 확률은 공준 1 에 의해 그때 저장조가 가질 수 있는 경우의
수에 비례한다:
\begin{equation}
P_i\;\propto\;W_R(E_\mathrm{tot}-E_i,\,N_\mathrm{tot}-N_i)
\;=\;\exp\!\Big[\frac{S_R(E_\mathrm{tot}-E_i,\,N_\mathrm{tot}-N_i)}{k_B}\Big].
\label{eq:sm-resv}
\end{equation}
저장조가 계보다 훨씬 크므로 $S_R$ 을 1차에서 절단해 전개하면(식~\eqref{eq:sm-fund} 의 두 도함수 대입)
\begin{equation}
S_R(E_\mathrm{tot}-E_i,\,N_\mathrm{tot}-N_i)
= S_R(E_\mathrm{tot},N_\mathrm{tot})-\frac{E_i}{T}+\frac{\mu N_i}{T}+\mathcal O(E_i^2,N_i^2).
\label{eq:sm-taylor}
\end{equation}
\textbf{(c) 중간식 --- Gibbs 인자와 canonical 특수형.} 식~\eqref{eq:sm-taylor} 을 \eqref{eq:sm-resv} 에 넣으면
상수 몫이 규격화로 빠지고
\begin{equation}
P_i\;\propto\;e^{-\beta(E_i-\mu N_i)},\qquad \beta\equiv\frac{1}{k_BT}
\label{eq:sm-gibbsfactor}
\end{equation}
--- 이 지수가 \emph{Gibbs 인자}다. 입자 교환이 없으면($N_i$ 고정) $e^{-\beta E_i}$ 의 \emph{Boltzmann 인자}로
줄고, 그 규격화 합이 canonical 분배함수(partition function) $Z=\sum_ie^{-\beta E_i}$, 그 로그가 Helmholtz
자유에너지 $F=-k_BT\ln Z$ 다($\langle E\rangle=\partial(\beta F)/\partial\beta$·$S=-\partial F/\partial T$ 가
열역학 관계와 일치함이 이 명명의 근거다). 이 $F$ 는 Faraday 상수 $F$ 와 기호만 같고 문맥으로 갈린다 ---
전자는 $\partial F/\partial T$ 처럼 \emph{미분당하는 자유에너지}, 후자는 $FV$·$RT/F$ 의 전하 환산
곱$\cdot$나눗셈 상수다. 식~\eqref{eq:sm-fund} 의 $\mu$ 를 $F$ 로 옮기면
$\mu=\partial F/\partial N|_{T,V}$ --- \emph{온도·부피를 고정한 채 입자 1개를 더 넣는 자유에너지 비용}이며,
이것이 본론 \S\ref{sec:center} 가 쓰는 거시 정의 $\mu=\partial G/\partial n|_{T,P}$ 와 같은 양임은
\S\ref{sec:sm-macro} 에서 닫는다.

\textbf{(d) 박스 --- grand canonical 분배함수.} 입자수까지 교환하는 앙상블이
\emph{대정준}(grand canonical) 앙상블이고, 그 규격화 합과 평균 입자수는
\begin{equation}
\boxed{\;\Xi=\sum_i e^{-\beta(E_i-\mu N_i)},\qquad
P_i=\frac{e^{-\beta(E_i-\mu N_i)}}{\Xi},\qquad
\langle N\rangle=k_BT\,\frac{\partial\ln\Xi}{\partial\mu}\;}
\label{eq:sm-gc}
\end{equation}
($\partial\ln\Xi/\partial\mu=\beta\sum_iN_iP_i=\beta\langle N\rangle$ 로 셋째 식이 첫째·둘째에서 나온다).
\begin{verifybox}
극한 검산 --- $\beta\to0$ 이면 $P_i\to$ 등확률(열적 규율 소실), $\beta\to\infty$ 면 $E_i-\mu N_i$ 최소 상태만
남는다. 삽입 전극이 정확히 이 설정이다 --- 격자의 자리들이 전해질을 통해 Li 저장조(상대극 Li 금속)와 입자를
교환하므로 $\mu$ 가 고정된 대정준 문제이고(그림~\ref{fig:sm-reservoir}), 다음 소절이 그 최소 단위(자리
하나)를 푼다.
\end{verifybox}

\begin{figure}[t]
\centering
\begin{tikzpicture}[font=\scriptsize,
  box/.style={draw,rounded corners=2pt,align=center,inner sep=4pt},
  ar/.style={{Latex[length=1.5mm]}-{Latex[length=1.5mm]},thick}]
\node[box,minimum width=30mm,minimum height=16mm] (sys) {system\\insertion sites\\microstate $i$: $(E_i,\,N_i)$};
\node[box,minimum width=46mm,minimum height=16mm,fill=black!6,right=20mm of sys] (res)
  {reservoir: $T,\ \mu$ fixed\\Li metal counter electrode\\$+$ electrolyte};
\draw[ar] ([yshift=3.2mm]sys.east) -- ([yshift=3.2mm]res.west) node[midway,above] {energy $E$};
\draw[ar] ([yshift=-3.2mm]sys.east) -- ([yshift=-3.2mm]res.west) node[midway,below] {particles $N$ (Li)};
\end{tikzpicture}
\caption{grand canonical 구도. 관심 계(삽입 자리들)는 상대극$\cdot$전해질이 이루는 큰 저장조와
에너지$\cdot$입자를 모두 교환하고, 저장조는 $T$ 와 $\mu$ 를 고정한다. 에너지만 열면 Boltzmann 인자, 입자까지
열면 Gibbs 인자(식~\eqref{eq:sm-gibbsfactor}) --- 지수의 자리가 하나 늘어날 뿐 논리는 같다.}
\label{fig:sm-reservoir}
\end{figure}

\subsection{단일 삽입 자리 --- 점유 통계와 그 요동}\label{sec:sm-site}
\textbf{(a) 출발 --- 두 가정과 미시상태의 온전한 기술.} 격자의 삽입 자리 \emph{하나}를 떼어 살펴본다.
자리는 전해질 너머의 Li 저장조와 입자를 교환하므로 식~\eqref{eq:sm-gc} 의 대정준 설정 그대로다. 유도에
앞서 두 가정을 명시한다.
\begin{itemize}[leftmargin=2.2em,itemsep=0.1em]
\item \textbf{가정 1 (hard-core 배타).} 한 자리에는 Li 이온이 $0$ 개 또는 $1$ 개만 들어간다. 근거는 양자
통계가 아니라 강한 단거리 반발이다 --- 자리의 부피와 이온 사이 정전$\cdot$입체 반발이 이중 점유의
에너지를 사실상 무한대로 올린다. 이 가정이 점유수 합을 $n\in\{0,1\}$ 로 자른다.
\item \textbf{가정 2 (자리 독립).} 이 소절에서는 자리의 에너지가 이웃 자리의 점유와 무관하다고 둔다.
이 가정은 \S\ref{sec:sm-mf} 에서 상호작용 파라미터 $\Omega$ 로 완화된다.
\end{itemize}
점유수 $n$ 만으로는 미시상태가 다 지정되지 않는다는 점이 중요하다. $n=1$ 인 자리에서 이온은 한 점에
얼어붙어 있는 것이 아니라 자리 우물 안에서 진동하는 연속 자유도(위치$\cdot$운동량)를 가진다. 따라서
미시상태의 목록은 ``빈 자리($n=0$, 에너지 $0$ 기준)'' 하나와, ``점유 자리($n=1$, 우물 바닥 에너지
$\varepsilon_0$)에 내부 배치 $\Gamma$ 가 붙은 연속족''이다.

\textbf{(b) 연산 --- 점유수 합 $\times$ 내부 자유도 적분.} 식~\eqref{eq:sm-gc} 의 대정준 합을 이 목록
그대로 쓰면, 점유수에 대한 합(가정 1 로 두 항)과 각 점유 상태의 내부 배치에 대한 적분으로 갈라진다:
\begin{equation}
\Xi_1\;=\;\sum_{n=0,1}e^{\beta\mu n}\,z_n
\;=\;1+q(T)\,e^{-\beta(\varepsilon_0-\mu)},
\qquad
q(T)\equiv\int e^{-\beta H_\mathrm{int}(\Gamma)}\,\frac{\dd\Gamma}{h^3}
\label{eq:partfn}
\end{equation}
--- $z_n$ 은 점유수 $n$ 인 상태들의 canonical 분배함수로, $z_0=1$(빈 자리는 상태 하나),
$z_1=e^{-\beta\varepsilon_0}q(T)$(우물 바닥 인자에 내부 자유도의 적분 $q(T)$ 가 곱해진다)이다.
$q(T)$ 는 점유된 이온의 진동 등 국소 자유도의 단일 자리 분배함수다\footnote{이 $q(T)$(단일 자리 내부 분배함수)는 N6--N9 의 정규화 용량 좌표 $q{=}Q/Q_\cell$ 와 다른 양이다 --- 전자는 Part 0 의 국소 열역학량, 후자는 곡선 좌표.}. 일반적인 정의는 상태
합 $q(T)=\mathrm{Tr}\,e^{-\beta H_\mathrm{int}}$ 이고 식~\eqref{eq:partfn} 의 위상공간 적분은 그 고전
극한이다 --- 자리 우물을 3차원 조화 퍼텐셜(진동수 $\omega_0$)로 근사하면 고전 극한에서
$q=(k_BT/\hbar\omega_0)^3$ 이고, 등방 가정(세 방향이 같은 $\omega_0$)을 풀어 데카르트 축마다 다른
진동수를 허용하는 일반형으로 넓히면 양자 상태 합으로는
$q=\prod_{i=1}^{3}[2\sinh(\beta\hbar\omega_i/2)]^{-1}$(축 지표 $i=1,2,3$ 마다 진동수 $\omega_i$)로 닫힌다. 이 전개는 국소화
흡착의 Langmuir 등온선 표준 유도와 동일하고 \cite{hill1960,fowler1939,mcquarrie1976}, 삽입 전극에의
적용 원형은 McKinnon--Haering \cite{mckinnon1983} 이다. 기호에 관하여 --- $\Xi_1$ 은 자리 1개의
\emph{대정준} 분배함수로, \S\ref{sec:sm-ensemble} 의 canonical $Z$ 와는 앙상블이 다르므로 기호부터
구분해 적는다(첨자 $1$ = 자리 1개). Chapter 2 의 단일 자리 분배함수(그 문건 식 (eq:Z1))도 같은 기호
$\Xi_1$ 로 통일한다.

\textbf{(c) 중간식 --- 평균 점유와 계수 $q(T)$ 의 흡수.} 점유 상태들의 확률 합이 곧 평균 점유다:
\begin{equation}
\langle n\rangle=0\cdot P_0+1\cdot P_1
=\frac{q(T)\,e^{-\beta(\varepsilon_0-\mu)}}{1+q(T)\,e^{-\beta(\varepsilon_0-\mu)}},
\label{eq:sm-occmid}
\end{equation}
그리고 식~\eqref{eq:sm-gc} 의 셋째 식으로 교차검증하면 $k_BT\,\partial\ln\Xi_1/\partial\mu$ 가 같은
결과를 준다. 여기서 적분이 남긴 계수 $q(T)$ 는 \emph{분자와 분모의 점유-상태 항 양쪽에 같은 인수로}
붙어 있다. 따라서 결과는 $q$ 와 $\varepsilon_0$ 을 따로 아는 데 의존하지 않고 한 조합에만 의존하며,
그 조합은 지수 안으로 흡수된다:
\begin{equation}
q(T)\,e^{-\beta\varepsilon_0}\;=\;e^{-\beta\tilde\varepsilon(T)},
\qquad
\tilde\varepsilon(T)\;\equiv\;\varepsilon_0-k_BT\ln q(T)
\label{eq:sm-epstilde}
\end{equation}
--- $\tilde\varepsilon$ 는 점유 자리의 \emph{자유에너지}(우물 바닥 에너지 $+$ 내부 자유도의 자유에너지
$-k_BT\ln q$)다. 자리 점유의 자유에너지 비용을 $\Delta\mu\equiv\tilde\varepsilon-\mu_\mathrm{Li}$ 로
줄여 쓰면 식~\eqref{eq:sm-occmid} 는 $e^{-\beta\Delta\mu}/(1+e^{-\beta\Delta\mu})$ 가 되어, 내부
자유도가 있어도 함수형은 두-상태 점유의 logistic 그대로 보존된다.

\textbf{(d) 박스 --- 점유율 $\theta$.} 분모·분자를 $e^{-\beta\Delta\mu}$ 로 나누면
\begin{equation}
\boxed{\;\theta\;\equiv\;\langle n\rangle\;=\;\frac{1}{1+e^{+\beta\,\Delta\mu}}\;,
\qquad \Delta\mu\equiv\tilde\varepsilon(T)-\mu_\mathrm{Li}\;.}
\label{eq:fermifn}
\end{equation}
\begin{verifybox}
극한 검산 --- $\beta\to\infty$($T\to0$)에서 $\Delta\mu\gtrless0$ 에 따라 $\theta\to0/1$ 의 계단
($T\to0$ 에서 $\tilde\varepsilon$ 는 점유 상태의 바닥 에너지로 수렴한다 --- 고전 $q$ 면 $k_BT\ln q\to0$
이라 $\varepsilon_0$, 양자 조화 $q$ 면 영점 에너지를 더한 $\varepsilon_0+\tfrac32\hbar\omega_0$ --- 어느
쪽이든 상수라 계단 구조는 동일),
$\beta\to0$($T\to\infty$)에서는 $\beta\Delta\mu\to-\ln q$ 이므로 내부 자유도가 없는 기준($q\equiv1$)에서
$\theta\to\tfrac12$(두 상태 등확률)이고, $q(T)$ 가 온도와 함께 증가하는 실제 조화 우물에서는 점유 상태의
위상공간 증가가 이겨 $\theta\to1$ 로 향한다; $\mu_\mathrm{Li}\to\pm\infty$ 에서
$\theta\to1/0$; 내부 자유도가 없는 극한 $q\to1$ 에서 $\tilde\varepsilon\to\varepsilon_0$ 으로 소박한
2-상태 결과를 회수한다. ★함수형 가드 --- 식~\eqref{eq:fermifn} 은 Fermi--Dirac 함수와 \emph{동형}이지만
물리량은 다르다: 공유하는 것은 ``한 자리에 입자 0 또는 1''이라는 배타 점유의 대수 구조뿐이고(가정 1),
여기의 입자는 전자가 아니라 Li 이온이며 배타성의 근거도 양자 통계가 아니라 자리 하나에 이온 하나라는
기하다.
\end{verifybox}

한 걸음 더 --- $n\in\{0,1\}$ 라 $n^2=n$ 이므로 점유의 요동(분산)이 닫힌 꼴로 나온다:
\begin{equation}
\mathrm{var}(n)=\langle n^2\rangle-\langle n\rangle^2=\theta(1-\theta),
\qquad
\frac{\partial\theta}{\partial(\beta\mu)}=\theta(1-\theta)
\label{eq:sm-flucres}
\end{equation}
(둘째 식은 식~\eqref{eq:fermifn} 의 $\beta\mu$ 미분). \emph{점유의 $\mu$-감도 $=$ 점유 요동} --- 요동--응답
(fluctuation--response) 관계의 최소 사례이고, 본론 평형 peak 의 종 $\xi_\eq(1-\xi_\eq)$(\S\ref{sec:eqpeak})의
통계적 정체가 바로 이 분산이다.

유효 자리 자유에너지 $\tilde\varepsilon(T)$ 가 하류에서 가는 곳은 두 갈래다. 첫째, 온도를 고정하고 보면
$\tilde\varepsilon$ 는 상수 하나이고, \S\ref{sec:sm-electro} 의 전기화학 결선에서 전이 중심 $U_j$ 로
흡수된다 --- 그 소절과 본론이 쓰는 몰당 기준 에너지 $\mu^0$ 는 정확히 $\tilde\varepsilon$ 의 몰 환산
$N_A\tilde\varepsilon$ 이다. 둘째, 온도를 미분하면 내부 자유도의 자유에너지 몫
$f_\mathrm{int}=-k_BT\ln q$ 가 자리당 엔트로피
\begin{equation}
s_\mathrm{int}\;=\;-\frac{\partial f_\mathrm{int}}{\partial T}
\;=\;k_B\,\frac{\partial\,(T\ln q)}{\partial T}
\label{eq:sm-sint}
\end{equation}
를 주는데, 이것이 Chapter 2 의 vibrational 엔트로피 사슬의 단일 자리 원형이다 --- 여기 자리 하나의
세 축 진동수 $\omega_i$($i=1,2,3$)를 격자 전체의 모든 진동 모드를 헤아리는 지표 $k$ 로 넓혀, 조화 모드의
$q_k=[1-e^{-\beta\hbar\omega_k}]^{-1}$ 를 식~\eqref{eq:sm-sint} 에 넣으면 그 문건의 모드당 엔트로피
$s_k=k_B[(1+\bar n_k)\ln(1+\bar n_k)-\bar n_k\ln\bar n_k]$(Bose--Einstein 들뜸수 $\bar n_k$)가 그대로
나온다. 기준점에 관한 사소한 주의로, 위 (b) 의 $[2\sinh(\beta\hbar\omega/2)]^{-1}$ 은 우물 바닥 기준
(영점 에너지를 $q$ 에 포함)이고 여기의 $[1-e^{-\beta\hbar\omega}]^{-1}$ 은 영점 에너지를
$\varepsilon_0$ 쪽에 포함시킨 기준이다 --- 두 규약의 $q$ 는 인자 $e^{-\beta\hbar\omega/2}$(영점 에너지
$\hbar\omega/2$ 의 Boltzmann 인자)만큼 다른데, 이 인자가 $T\ln q$ 에 남기는 기여는 온도 무관 상수
$-\hbar\omega/(2k_B)$ 이므로 온도 미분, 곧 엔트로피는 동일하다. 따라서 전이 중심의 온도 기울기
$\dd U_j/\dd T=\Delta S_{\rxn,j}/F$(\S\ref{sec:center})에서 진동
엔트로피가 차지하는 몫의 미시적 기원이 바로 $\tilde\varepsilon(T)$ 의 온도 의존, 곧 $q(T)$ 다.

\subsection{$M$ 개의 독립 자리 --- lattice gas 와 섞임 엔트로피}\label{sec:sm-lattice}
\textbf{(a) 출발 --- 동등·독립 자리의 집합.} 결정 안에 동등한 삽입 자리가 $M$ 개 있고 서로 상호작용하지
않는다고 하자 --- 이 모형이 \emph{격자기체}(lattice gas)의 이상 극한이다. 미시상태는 점유 배열
$(n_1,\dots,n_M)$($n_k\in\{0,1\}$, 가정 1)에 각 점유 자리의 내부 배치가 붙은 것이고, $N=\sum_kn_k$ 다.

\textbf{(b) 연산 --- 분배함수의 인수분해.} 자리들이 독립(가정 2)이라 대정준 합이 자리별 곱으로 갈라지고,
각 인수는 내부 자유도 적분까지 마친 단일 자리 결과~\eqref{eq:partfn} 그대로다
($\Delta\mu=\tilde\varepsilon-\mu$ 로 흡수된 표기):
\begin{equation}
\Xi_M=\sum_{n_1=0,1}\!\cdots\!\sum_{n_M=0,1}\prod_{k=1}^{M}e^{-\beta\,\Delta\mu\,n_k}
=\prod_{k=1}^{M}\Big[\sum_{n_k=0,1}e^{-\beta\,\Delta\mu\,n_k}\Big]=\Xi_1^{\,M},
\label{eq:sm-factor}
\end{equation}
따라서 식~\eqref{eq:sm-gc} 의 셋째 식으로 $\langle N\rangle=k_BT\,\partial(M\ln\Xi_1)/\partial\mu=M\theta$ ---
거시 점유율이 단일 자리 점유 확률 \eqref{eq:fermifn} 와 같다.

\textbf{(c) 중간식 --- 같은 결과의 셈법 경로(교차검증).} 이번엔 입자수 $N$ 을 고정하고 배치를 센다.
$M$ 자리에서 $N$ 자리를 고르는 경우의 수 $W=M!/[N!(M-N)!]$ 에 공준 2(식~\eqref{eq:sm-S})와 Stirling 근사
$\ln m!\simeq m\ln m-m$ 을 적용하면($\theta=N/M$)
\begin{equation}
S_\mathrm{mix}=k_B\ln W\simeq-k_BM\big[\theta\ln\theta+(1-\theta)\ln(1-\theta)\big],
\label{eq:sm-Smix}
\end{equation}
자유에너지 $F(N)=N\tilde\varepsilon-TS_\mathrm{mix}$(점유 자리당 자유에너지는 내부 자유도 몫까지 포함한
$\tilde\varepsilon$, 식~\eqref{eq:sm-epstilde})를 $N$ 으로 미분하면
\begin{equation}
\mu=\frac{\partial F}{\partial N}
=\tilde\varepsilon+k_BT\ln\frac{\theta}{1-\theta}.
\label{eq:sm-mucount}
\end{equation}
한편 식~\eqref{eq:fermifn} 를 $\mu$ 에 대해 뒤집어도($e^{\beta\Delta\mu}=(1-\theta)/\theta$) 같은
식~\eqref{eq:sm-mucount} 가 나온다 --- 대정준 경로(자리 하나의 확률)와 셈법 경로(배치 수의 로그)가 같은
$\mu(\theta)$ 에 닿는 것이 앙상블 동등성이며, 본론 \S\ref{sec:dist} 가 말하는 ``kinetic·thermo 두 경로가 한
분포의 두 얼굴''의 평형 쪽 절반이 바로 이것이다.

\textbf{(d) 박스 --- 몰 단위 이상 lattice gas 화학퍼텐셜.} 자리당 $k_B$·$\tilde\varepsilon$ 을 몰당
$R=N_Ak_B$·$\mu^0\equiv N_A\tilde\varepsilon$ 으로 환산하면
\begin{equation}
\boxed{\;\mu(\theta)=\mu^0+RT\ln\frac{\theta}{1-\theta}\qquad(\text{이상 lattice gas})\;.}
\label{eq:sm-muideal}
\end{equation}
\begin{verifybox}
부호·극한 검산 --- $\theta\to0$ 에서 $\mu\to-\infty$(거의 빈 격자에 입자를 넣는 일은 섞임 엔트로피 이득 때문에
얼마든지 값싸다), $\theta\to1$ 에서 $+\infty$(마지막 빈자리는 무한히 비싸다), $\theta=\tfrac12$ 에서
$\mu=\mu^0$·자리당 $S_\mathrm{mix}$ 최대 $k_B\ln2$. 식~\eqref{eq:sm-muideal} 는 다음 소절 $\mu(\theta)$ 의
$\Omega=0$ 특수형이고, 상호작용 몫 $\Omega(1-2\theta)$ 를 다음 소절이 더한다.
\end{verifybox}

% 자산: [A-015] [A-016] [A-017] [A-018] [A-019] [A-020] [A-021] [A-022] [A-023] [A-024] [A-025] [A-026] [A-027] [A-028] [A-029] [A-030] [A-031] [A-032] [A-033] [A-034] [A-035] [A-036] [A-037] [A-038] [A-039] [A-040] [A-041] [A-042] [A-043] [A-044] [A-045] [A-046]
